\section{Accelerator / DPU / SmartNIC confidential offload}

\begin{frame}[t,shrink=6]{Accelerator / DPU / SmartNIC confidential offload：方向开场}
\DeckKicker{方向开场}
\SlideMeta{证据等级}{方向综述}{来源状态：公开材料综合}
{\large\bfseries\color{Ink} Accelerator / DPU / SmartNIC confidential offload 的核心是先界定保护对象、管理边界和证据强度，再用三篇主讲材料串起技术演进。\par}\vspace{1.2mm}
\begin{columns}[T,totalwidth=\textwidth]
  \begin{column}{0.45\textwidth}
    \fcolorbox{Rule}{Paper}{%
  \begin{minipage}[t]{0.97\linewidth}
    \vspace{1.1mm}
    \hspace{1.4mm}\begin{minipage}{0.92\linewidth}
      \PanelTitle{内容说明}
      \scriptsize \begin{itemize}
  \item 从 rack-level accelerator TEE 走向 cloud-scale DSA/NPU/GPU pool 与 Arm CCA accelerator。
  \item 固定三篇主讲材料；`sang2025portal`、`deng2022strongbox`、`nvidia\_bluefield\_operation\_2025` 作为辅助对比，不计入三篇主讲。
  \item 先讲方向痛点和设计轴，再逐篇说明贡献、限制、证据强度和可落地条件。
  \item 对规范、Survey、draft、vendor evidence 保持显式标注，避免把背景材料写成一手系统证明。
\end{itemize}
    \end{minipage}
    \vspace{1mm}
  \end{minipage}}
  \end{column}
  \begin{column}{0.49\textwidth}
    \fcolorbox{Rule}{Panel}{%
  \begin{minipage}[t]{0.97\linewidth}
    \vspace{1.1mm}
    \hspace{1.4mm}\begin{minipage}{0.92\linewidth}
      \PanelTitle{三篇主讲材料的分工}
      \scriptsize {\tiny
\begin{tabularx}{\linewidth}{@{}YYY@{}}
\toprule
\textbf{论文} & \textbf{定位} & \textbf{证据边界} \\
\midrule
HETEE: rack-scale heterogeneous TEE & 开创/基础入口 & E1 / 基础证据 peer-reviewed \\
CloudScale confidential computing with heterogeneous devices & 代表性改进 & E1 / 同行评审改进证据 \\
CAGE confidential accelerator workflowow for Arm CCA & 当前 SOTA/补充边界 & E1 / 同行评审改进证据 \\
\bottomrule
\end{tabularx}
}
    \end{minipage}
    \vspace{1mm}
  \end{minipage}}
  \end{column}
\end{columns}
\SourceFootnote{来源：论文原文、官方规范或公开材料｜证据等级：见本页 badge}
\end{frame}


\begin{frame}[t,shrink=6]{HETEE: rack-scale heterogeneous TEE：内容摘要}
\DeckKicker{内容摘要}
\SlideMeta{证据等级}{E1 / 基础证据 peer-reviewed}{来源状态：本地 PDF 已核验}
{\large\bfseries\color{Ink} HETEE 用 rack-level security controller 与 PCIe fabric 动态隔离 GPU/FPGA/TPU 等。\par}\vspace{1.2mm}
\begin{columns}[T,totalwidth=\textwidth]
  \begin{column}{0.45\textwidth}
    \fcolorbox{Rule}{Paper}{%
  \begin{minipage}[t]{0.97\linewidth}
    \vspace{1.1mm}
    \hspace{1.4mm}\begin{minipage}{0.92\linewidth}
      \PanelTitle{内容说明}
      \scriptsize \begin{itemize}
  \item HETEE 用 rack-level security controller 与 PCIe fabric 动态隔离 GPU/FPGA/TPU 等 accelerator
  \item HETEE is a foundational peer-reviewed work for heterogeneous CPU-accelerator TEEs.
\end{itemize}
    \end{minipage}
    \vspace{1mm}
  \end{minipage}}
  \end{column}
  \begin{column}{0.49\textwidth}
    \fcolorbox{Rule}{Panel}{%
  \begin{minipage}[t]{0.97\linewidth}
    \vspace{1.1mm}
    \hspace{1.4mm}\begin{minipage}{0.92\linewidth}
      \PanelTitle{证据定位}
      \scriptsize {\tiny
\begin{tabularx}{\linewidth}{@{}YY@{}}
\toprule
\textbf{字段} & \textbf{内容} \\
\midrule
论文定位 & 开创/基础入口 \\
材料类型 & 系统论文 \\
证据等级 & E1 / 基础证据 peer-reviewed \\
来源状态 & 本地 PDF 已核验 \\
\bottomrule
\end{tabularx}
}
    \end{minipage}
    \vspace{1mm}
  \end{minipage}}
  \end{column}
\end{columns}
\SourceFootnote{来源：HETEE: rack-scale heterogeneous TEE｜证据等级：E1 / 基础证据 peer-reviewed}
\end{frame}


\begin{frame}[t,shrink=6]{HETEE: rack-scale heterogeneous TEE：研究背景}
\DeckKicker{研究背景}
\SlideMeta{证据等级}{E1 / 基础证据 peer-reviewed}{来源状态：本地 PDF 已核验}
{\large\bfseries\color{Ink} CPU TEE 不能天然保护 accelerator offload path，敏感数据进入 GPU、driver、PCIe fabric 或 remote resource pool 后仍可能暴露\par}\vspace{1.2mm}
\begin{columns}[T,totalwidth=\textwidth]
  \begin{column}{0.45\textwidth}
    \fcolorbox{Rule}{Paper}{%
  \begin{minipage}[t]{0.97\linewidth}
    \vspace{1.1mm}
    \hspace{1.4mm}\begin{minipage}{0.92\linewidth}
      \PanelTitle{内容说明}
      \scriptsize \begin{itemize}
  \item CPU TEE 不能天然保护 accelerator offload path，敏感数据进入 GPU、driver、PCIe fabric 或 remote resource...
\end{itemize}
    \end{minipage}
    \vspace{1mm}
  \end{minipage}}
  \end{column}
  \begin{column}{0.49\textwidth}
    \fcolorbox{Rule}{Panel}{%
  \begin{minipage}[t]{0.97\linewidth}
    \vspace{1.1mm}
    \hspace{1.4mm}\begin{minipage}{0.92\linewidth}
      \PanelTitle{问题边界}
      \scriptsize {\tiny
\begin{tabularx}{\linewidth}{@{}YY@{}}
\toprule
\textbf{维度} & \textbf{说明} \\
\midrule
保护对象 & Accelerator / DPU / SmartNIC confidential offload \\
传统瓶颈 & 从 rack-level accelerator TEE 走向 cloud-scale DSA/NPU/GPU pool 与 Arm... \\
本文入口 & HETEE: rack-scale heterogeneous TEE \\
证据限制 & E1 / 基础证据 peer-reviewed \\
\bottomrule
\end{tabularx}
}
    \end{minipage}
    \vspace{1mm}
  \end{minipage}}
  \end{column}
\end{columns}
\SourceFootnote{来源：HETEE: rack-scale heterogeneous TEE｜证据等级：E1 / 基础证据 peer-reviewed}
\end{frame}


\begin{frame}[t,shrink=6]{HETEE: rack-scale heterogeneous TEE：解决方案}
\DeckKicker{解决方案}
\SlideMeta{证据等级}{E1 / 基础证据 peer-reviewed}{来源状态：本地 PDF 已核验}
{\large\bfseries\color{Ink} Security Controller 作为 thin TCB 管理 accelerator assignment、加密认证、remote attestation、secure cleanup 和资源回收\par}\vspace{1.2mm}
\begin{columns}[T,totalwidth=\textwidth]
  \begin{column}{0.45\textwidth}
    \fcolorbox{Rule}{Paper}{%
  \begin{minipage}[t]{0.97\linewidth}
    \vspace{1.1mm}
    \hspace{1.4mm}\begin{minipage}{0.92\linewidth}
      \PanelTitle{内容说明}
      \scriptsize \begin{itemize}
  \item Security Controller 作为 thin TCB 管理 accelerator assignment、加密认证、remote attestation、secure...
\end{itemize}
    \end{minipage}
    \vspace{1mm}
  \end{minipage}}
  \end{column}
  \begin{column}{0.49\textwidth}
    \fcolorbox{Rule}{Panel}{%
  \begin{minipage}[t]{0.97\linewidth}
    \vspace{1.1mm}
    \hspace{1.4mm}\begin{minipage}{0.92\linewidth}
      \PanelTitle{核心设计拆解}
      \scriptsize \begin{itemize}
  \item 01. Rack-level Security Controller
  \item 02. PCIe ExpressFabric accelerator assignment
  \item 03. Remote attestation and secure cleanup lifecycle
\end{itemize}
    \end{minipage}
    \vspace{1mm}
  \end{minipage}}
  \end{column}
\end{columns}
\SourceFootnote{来源：HETEE: rack-scale heterogeneous TEE｜证据等级：E1 / 基础证据 peer-reviewed}
\end{frame}


\begin{frame}[t,shrink=6]{HETEE: rack-scale heterogeneous TEE：实验结果}
\DeckKicker{实验结果}
\SlideMeta{证据等级}{E1 / 基础证据 peer-reviewed}{来源状态：本地 PDF 已核验}
{\large\bfseries\color{Ink} 真实硬件原型\par}\vspace{1.2mm}
\begin{columns}[T,totalwidth=\textwidth]
  \begin{column}{0.45\textwidth}
    \fcolorbox{Rule}{Paper}{%
  \begin{minipage}[t]{0.97\linewidth}
    \vspace{1.1mm}
    \hspace{1.4mm}\begin{minipage}{0.92\linewidth}
      \PanelTitle{内容说明}
      \scriptsize \begin{itemize}
  \item 真实硬件原型
  \item 论文报告 ResNet152 inference 最大 throughput overhead 2.17\%，training 最大 overhead 0.95\%
\end{itemize}
    \end{minipage}
    \vspace{1mm}
  \end{minipage}}
  \end{column}
  \begin{column}{0.49\textwidth}
    \fcolorbox{Rule}{Panel}{%
  \begin{minipage}[t]{0.97\linewidth}
    \vspace{1.1mm}
    \hspace{1.4mm}\begin{minipage}{0.92\linewidth}
      \PanelTitle{实验/证据边界}
      \scriptsize {\tiny
\begin{tabularx}{\linewidth}{@{}YYY@{}}
\toprule
\textbf{对象} & \textbf{可支撑} & \textbf{边界} \\
\midrule
数据/来源 & HETEE: rack-scale heterogeneous TEE & 本地 PDF 已核验 \\
Baseline/对照 & 以论文或规范原文为准 & 不跨材料外推 \\
指标/对象 & 只使用当前来源可证明的信息 & E1 / 基础证据 peer-reviewed \\
使用边界 & 可进入本方向演进图 & 不能替代更强证据 \\
\bottomrule
\end{tabularx}
}
    \end{minipage}
    \vspace{1mm}
  \end{minipage}}
  \end{column}
\end{columns}
\SourceFootnote{来源：HETEE: rack-scale heterogeneous TEE｜证据等级：E1 / 基础证据 peer-reviewed}
\end{frame}


\begin{frame}[t,shrink=6]{HETEE: rack-scale heterogeneous TEE：文章评价}
\DeckKicker{文章评价}
\SlideMeta{证据等级}{E1 / 基础证据 peer-reviewed}{来源状态：本地 PDF 已核验}
{\large\bfseries\color{Ink} accelerator confidential offload 的 foundational baseline，说明 CPU TEE 边界可以被。\par}\vspace{1.2mm}
\begin{columns}[T,totalwidth=\textwidth]
  \begin{column}{0.45\textwidth}
    \fcolorbox{Rule}{Paper}{%
  \begin{minipage}[t]{0.97\linewidth}
    \vspace{1.1mm}
    \hspace{1.4mm}\begin{minipage}{0.92\linewidth}
      \PanelTitle{内容说明}
      \scriptsize \begin{itemize}
  \item 设计优势：accelerator confidential offload 的 foundational baseline，说明 CPU TEE 边界可以被。
  \item 局限性：早于 SPDM/TDISP/CoVE-IO，不能当作 modern trusted I/O 标准或 vendor DPU attestation proof。
  \item 商业化潜力：适合启发 DPU/SmartNIC security-controller 设计；风险在集中 TCB、设备固件可信和 production lifecycle 复杂度。
  \item 证据边界：E1 / Foundational peer-reviewed；本地 PDF 已核验。
\end{itemize}
    \end{minipage}
    \vspace{1mm}
  \end{minipage}}
  \end{column}
  \begin{column}{0.49\textwidth}
    \fcolorbox{Rule}{Panel}{%
  \begin{minipage}[t]{0.97\linewidth}
    \vspace{1.1mm}
    \hspace{1.4mm}\begin{minipage}{0.92\linewidth}
      \PanelTitle{文章评价}
      \scriptsize {\tiny
\begin{tabularx}{\linewidth}{@{}YY@{}}
\toprule
\textbf{维度} & \textbf{判断} \\
\midrule
设计优势 & accelerator confidential offload 的 foundational baseline，说明 CPU TEE 边界可以被。 \\
主要局限 & 早于 SPDM/TDISP/CoVE-IO，不能当作 modern trusted I/O 标准或 vendor DPU attestation proof。 \\
商业落地 & 适合启发 DPU/SmartNIC security-controller 设计；风险在集中 TCB、设备固件可信和 production... \\
讲稿定位 & 开创/基础入口; 证据强度 E1 \\
\bottomrule
\end{tabularx}
}
    \end{minipage}
    \vspace{1mm}
  \end{minipage}}
  \end{column}
\end{columns}
\SourceFootnote{来源：HETEE: rack-scale heterogeneous TEE｜证据等级：E1 / 基础证据 peer-reviewed}
\end{frame}


\begin{frame}[t,shrink=6]{CloudScale confidential computing with heterogeneous devices：内容摘要}
\DeckKicker{内容摘要}
\SlideMeta{证据等级}{E1 / 同行评审改进证据}{来源状态：本地 PDF 已核验}
{\large\bfseries\color{Ink} CloudScale 用 cloud-scale Security Controller 桥接 TEE-enabled nodes 与 legacy non-TEE DSA/NPU/GPU nodes\par}\vspace{1.2mm}
\begin{columns}[T,totalwidth=\textwidth]
  \begin{column}{0.45\textwidth}
    \fcolorbox{Rule}{Paper}{%
  \begin{minipage}[t]{0.97\linewidth}
    \vspace{1.1mm}
    \hspace{1.4mm}\begin{minipage}{0.92\linewidth}
      \PanelTitle{内容说明}
      \scriptsize \begin{itemize}
  \item CloudScale 用 cloud-scale Security Controller 桥接 TEE-enabled nodes 与 legacy non-TEE...
  \item Cloud-scale acceleration work is a peer-reviewed 2024 SOTA item for confidential...
\end{itemize}
    \end{minipage}
    \vspace{1mm}
  \end{minipage}}
  \end{column}
  \begin{column}{0.49\textwidth}
    \fcolorbox{Rule}{Panel}{%
  \begin{minipage}[t]{0.97\linewidth}
    \vspace{1.1mm}
    \hspace{1.4mm}\begin{minipage}{0.92\linewidth}
      \PanelTitle{证据定位}
      \scriptsize {\tiny
\begin{tabularx}{\linewidth}{@{}YY@{}}
\toprule
\textbf{字段} & \textbf{内容} \\
\midrule
论文定位 & 代表性改进 \\
材料类型 & 系统论文 \\
证据等级 & E1 / 同行评审改进证据 \\
来源状态 & 本地 PDF 已核验 \\
\bottomrule
\end{tabularx}
}
    \end{minipage}
    \vspace{1mm}
  \end{minipage}}
  \end{column}
\end{columns}
\SourceFootnote{来源：CloudScale confidential computing with heterogeneous devices｜证据等级：E1 / 同行评审改进证据}
\end{frame}


\begin{frame}[t,shrink=6]{CloudScale confidential computing with heterogeneous devices：研究背景}
\DeckKicker{研究背景}
\SlideMeta{证据等级}{E1 / 同行评审改进证据}{来源状态：本地 PDF 已核验}
{\large\bfseries\color{Ink} 数据中心中多数 DSA 节点没有 TEE 能力\par}\vspace{1.2mm}
\begin{columns}[T,totalwidth=\textwidth]
  \begin{column}{0.45\textwidth}
    \fcolorbox{Rule}{Paper}{%
  \begin{minipage}[t]{0.97\linewidth}
    \vspace{1.1mm}
    \hspace{1.4mm}\begin{minipage}{0.92\linewidth}
      \PanelTitle{内容说明}
      \scriptsize \begin{itemize}
  \item 数据中心中多数 DSA 节点没有 TEE 能力
  \item 用户要在 accelerator 性能与敏感数据保护之间取舍
\end{itemize}
    \end{minipage}
    \vspace{1mm}
  \end{minipage}}
  \end{column}
  \begin{column}{0.49\textwidth}
    \fcolorbox{Rule}{Panel}{%
  \begin{minipage}[t]{0.97\linewidth}
    \vspace{1.1mm}
    \hspace{1.4mm}\begin{minipage}{0.92\linewidth}
      \PanelTitle{问题边界}
      \scriptsize {\tiny
\begin{tabularx}{\linewidth}{@{}YY@{}}
\toprule
\textbf{维度} & \textbf{说明} \\
\midrule
保护对象 & Accelerator / DPU / SmartNIC confidential offload \\
传统瓶颈 & 从 rack-level accelerator TEE 走向 cloud-scale DSA/NPU/GPU pool 与 Arm... \\
本文入口 & CloudScale confidential computing with heterogeneous devices \\
证据限制 & E1 / 同行评审改进证据 \\
\bottomrule
\end{tabularx}
}
    \end{minipage}
    \vspace{1mm}
  \end{minipage}}
  \end{column}
\end{columns}
\SourceFootnote{来源：CloudScale confidential computing with heterogeneous devices｜证据等级：E1 / 同行评审改进证据}
\end{frame}


\begin{frame}[t,shrink=6]{CloudScale confidential computing with heterogeneous devices：解决方案}
\DeckKicker{解决方案}
\SlideMeta{证据等级}{E1 / 同行评审改进证据}{来源状态：本地 PDF 已核验}
{\large\bfseries\color{Ink} SC 作为 TEE proxy 执行 access control、attestation、key exchange 和 encrypted/authenticated data path\par}\vspace{1.2mm}
\begin{columns}[T,totalwidth=\textwidth]
  \begin{column}{0.45\textwidth}
    \fcolorbox{Rule}{Paper}{%
  \begin{minipage}[t]{0.97\linewidth}
    \vspace{1.1mm}
    \hspace{1.4mm}\begin{minipage}{0.92\linewidth}
      \PanelTitle{内容说明}
      \scriptsize \begin{itemize}
  \item SC 作为 TEE proxy 执行 access control、attestation、key exchange 和 encrypted/authenticated data...
\end{itemize}
    \end{minipage}
    \vspace{1mm}
  \end{minipage}}
  \end{column}
  \begin{column}{0.49\textwidth}
    \fcolorbox{Rule}{Panel}{%
  \begin{minipage}[t]{0.97\linewidth}
    \vspace{1.1mm}
    \hspace{1.4mm}\begin{minipage}{0.92\linewidth}
      \PanelTitle{核心设计拆解}
      \scriptsize \begin{itemize}
  \item 01. Cloud-scale Security Controller
  \item 02. TEE / non-TEE DSA pool attestation
  \item 03. Encrypted and authenticated offload path
\end{itemize}
    \end{minipage}
    \vspace{1mm}
  \end{minipage}}
  \end{column}
\end{columns}
\SourceFootnote{来源：CloudScale confidential computing with heterogeneous devices｜证据等级：E1 / 同行评审改进证据}
\end{frame}


\begin{frame}[t,shrink=6]{CloudScale confidential computing with heterogeneous devices：实验结果}
\DeckKicker{实验结果}
\SlideMeta{证据等级}{E1 / 同行评审改进证据}{来源状态：本地 PDF 已核验}
{\large\bfseries\color{Ink} 论文报告 AI、Redis、file-system workloads 平均 1.5-5\% overhead，并扩展到 2236 concurrent NPUs running CNNs\par}\vspace{1.2mm}
\begin{columns}[T,totalwidth=\textwidth]
  \begin{column}{0.45\textwidth}
    \fcolorbox{Rule}{Paper}{%
  \begin{minipage}[t]{0.97\linewidth}
    \vspace{1.1mm}
    \hspace{1.4mm}\begin{minipage}{0.92\linewidth}
      \PanelTitle{内容说明}
      \scriptsize \begin{itemize}
  \item 论文报告 AI、Redis、file-system workloads 平均 1.5-5\% overhead，并扩展到 2236 concurrent NPUs running...
\end{itemize}
    \end{minipage}
    \vspace{1mm}
  \end{minipage}}
  \end{column}
  \begin{column}{0.49\textwidth}
    \fcolorbox{Rule}{Panel}{%
  \begin{minipage}[t]{0.97\linewidth}
    \vspace{1.1mm}
    \hspace{1.4mm}\begin{minipage}{0.92\linewidth}
      \PanelTitle{实验/证据边界}
      \scriptsize {\tiny
\begin{tabularx}{\linewidth}{@{}YYY@{}}
\toprule
\textbf{对象} & \textbf{可支撑} & \textbf{边界} \\
\midrule
数据/来源 & CloudScale confidential computing with heterogeneous devices & 本地 PDF 已核验 \\
Baseline/对照 & 以论文或规范原文为准 & 不跨材料外推 \\
指标/对象 & 只使用当前来源可证明的信息 & E1 / 同行评审改进证据 \\
使用边界 & 可进入本方向演进图 & 不能替代更强证据 \\
\bottomrule
\end{tabularx}
}
    \end{minipage}
    \vspace{1mm}
  \end{minipage}}
  \end{column}
\end{columns}
\SourceFootnote{来源：CloudScale confidential computing with heterogeneous devices｜证据等级：E1 / 同行评审改进证据}
\end{frame}


\begin{frame}[t,shrink=6]{CloudScale confidential computing with heterogeneous devices：文章评价}
\DeckKicker{文章评价}
\SlideMeta{证据等级}{E1 / 同行评审改进证据}{来源状态：本地 PDF 已核验}
{\large\bfseries\color{Ink} 把 accelerator TEE 从单机扩展到 cloud-scale heterogeneous DSA pool\par}\vspace{1.2mm}
\begin{columns}[T,totalwidth=\textwidth]
  \begin{column}{0.45\textwidth}
    \fcolorbox{Rule}{Paper}{%
  \begin{minipage}[t]{0.97\linewidth}
    \vspace{1.1mm}
    \hspace{1.4mm}\begin{minipage}{0.92\linewidth}
      \PanelTitle{内容说明}
      \scriptsize \begin{itemize}
  \item 设计优势：把 accelerator TEE 从单机扩展到 cloud-scale heterogeneous DSA pool。
  \item 局限性：SC 正确性和供应链可信是硬假设；不替代 SPDM/TDISP/PCIe IDE/IOMMU 或 vendor DPU RoT。
  \item 商业化潜力：贴近云 DPU/SmartNIC/NPU resource-pool 问题；风险在多租户编排、热插拔、failure isolation 和 production...
  \item 证据边界：E1 / Peer-reviewed SOTA；本地 PDF 已核验。
\end{itemize}
    \end{minipage}
    \vspace{1mm}
  \end{minipage}}
  \end{column}
  \begin{column}{0.49\textwidth}
    \fcolorbox{Rule}{Panel}{%
  \begin{minipage}[t]{0.97\linewidth}
    \vspace{1.1mm}
    \hspace{1.4mm}\begin{minipage}{0.92\linewidth}
      \PanelTitle{文章评价}
      \scriptsize {\tiny
\begin{tabularx}{\linewidth}{@{}YY@{}}
\toprule
\textbf{维度} & \textbf{判断} \\
\midrule
设计优势 & 把 accelerator TEE 从单机扩展到 cloud-scale heterogeneous DSA pool。 \\
主要局限 & SC 正确性和供应链可信是硬假设；不替代 SPDM/TDISP/PCIe IDE/IOMMU 或 vendor DPU RoT。 \\
商业落地 & 贴近云 DPU/SmartNIC/NPU resource-pool 问题；风险在多租户编排、热插拔、failure isolation 和。 \\
讲稿定位 & 代表性改进; 证据强度 E1 \\
\bottomrule
\end{tabularx}
}
    \end{minipage}
    \vspace{1mm}
  \end{minipage}}
  \end{column}
\end{columns}
\SourceFootnote{来源：CloudScale confidential computing with heterogeneous devices｜证据等级：E1 / 同行评审改进证据}
\end{frame}


\begin{frame}[t,shrink=6]{CAGE confidential accelerator workflowow for Arm CCA：内容摘要}
\DeckKicker{内容摘要}
\SlideMeta{证据等级}{E1 / 同行评审改进证据}{来源状态：本地 PDF 已核验}
{\large\bfseries\color{Ink} CAGE 用 Arm CCA GPC/GPT 与 Monitor-side shadow task 保护 Realm 使用 GPU/FPGA accelerator 的 workflowow\par}\vspace{1.2mm}
\begin{columns}[T,totalwidth=\textwidth]
  \begin{column}{0.45\textwidth}
    \fcolorbox{Rule}{Paper}{%
  \begin{minipage}[t]{0.97\linewidth}
    \vspace{1.1mm}
    \hspace{1.4mm}\begin{minipage}{0.92\linewidth}
      \PanelTitle{内容说明}
      \scriptsize \begin{itemize}
  \item CAGE 用 Arm CCA GPC/GPT 与 Monitor-side shadow task 保护 Realm 使用 GPU/FPGA accelerator 的。
  \item CAGE is a recent peer-reviewed SOTA item for GPU accelerator enclaves.
\end{itemize}
    \end{minipage}
    \vspace{1mm}
  \end{minipage}}
  \end{column}
  \begin{column}{0.49\textwidth}
    \fcolorbox{Rule}{Panel}{%
  \begin{minipage}[t]{0.97\linewidth}
    \vspace{1.1mm}
    \hspace{1.4mm}\begin{minipage}{0.92\linewidth}
      \PanelTitle{证据定位}
      \scriptsize {\tiny
\begin{tabularx}{\linewidth}{@{}YY@{}}
\toprule
\textbf{字段} & \textbf{内容} \\
\midrule
论文定位 & 当前 SOTA/补充边界 \\
材料类型 & 系统论文 \\
证据等级 & E1 / 同行评审改进证据 \\
来源状态 & 本地 PDF 已核验 \\
\bottomrule
\end{tabularx}
}
    \end{minipage}
    \vspace{1mm}
  \end{minipage}}
  \end{column}
\end{columns}
\SourceFootnote{来源：CAGE confidential accelerator workflowow for Arm CCA｜证据等级：E1 / 同行评审改进证据}
\end{frame}


\begin{frame}[t,shrink=6]{CAGE confidential accelerator workflowow for Arm CCA：研究背景}
\DeckKicker{研究背景}
\SlideMeta{证据等级}{E1 / 同行评审改进证据}{来源状态：本地 PDF 已核验}
{\large\bfseries\color{Ink} Arm CCA 保护 Realm memory，但 untrusted accelerator driver/runtime 会处理 code、metadata、MMIO、DMA buffer 和 completion...\par}\vspace{1.2mm}
\begin{columns}[T,totalwidth=\textwidth]
  \begin{column}{0.45\textwidth}
    \fcolorbox{Rule}{Paper}{%
  \begin{minipage}[t]{0.97\linewidth}
    \vspace{1.1mm}
    \hspace{1.4mm}\begin{minipage}{0.92\linewidth}
      \PanelTitle{内容说明}
      \scriptsize \begin{itemize}
  \item Arm CCA 保护 Realm memory，但 untrusted accelerator driver/runtime 会处理 code、metadata、MMIO、DMA...
\end{itemize}
    \end{minipage}
    \vspace{1mm}
  \end{minipage}}
  \end{column}
  \begin{column}{0.49\textwidth}
    \fcolorbox{Rule}{Panel}{%
  \begin{minipage}[t]{0.97\linewidth}
    \vspace{1.1mm}
    \hspace{1.4mm}\begin{minipage}{0.92\linewidth}
      \PanelTitle{问题边界}
      \scriptsize {\tiny
\begin{tabularx}{\linewidth}{@{}YY@{}}
\toprule
\textbf{维度} & \textbf{说明} \\
\midrule
保护对象 & Accelerator / DPU / SmartNIC confidential offload \\
传统瓶颈 & 从 rack-level accelerator TEE 走向 cloud-scale DSA/NPU/GPU pool 与 Arm... \\
本文入口 & CAGE confidential accelerator workflowow for Arm CCA \\
证据限制 & E1 / 同行评审改进证据 \\
\bottomrule
\end{tabularx}
}
    \end{minipage}
    \vspace{1mm}
  \end{minipage}}
  \end{column}
\end{columns}
\SourceFootnote{来源：CAGE confidential accelerator workflowow for Arm CCA｜证据等级：E1 / 同行评审改进证据}
\end{frame}


\begin{frame}[t,shrink=6]{CAGE confidential accelerator workflowow for Arm CCA：解决方案}
\DeckKicker{解决方案}
\SlideMeta{证据等级}{E1 / 同行评审改进证据}{来源状态：本地 PDF 已核验}
{\large\bfseries\color{Ink} 用 shadow task、Monitor verification、GPC/GPT protection 和 accelerator-specific cleanup，把真实 accelerator task 从。\par}\vspace{1.2mm}
\begin{columns}[T,totalwidth=\textwidth]
  \begin{column}{0.45\textwidth}
    \fcolorbox{Rule}{Paper}{%
  \begin{minipage}[t]{0.97\linewidth}
    \vspace{1.1mm}
    \hspace{1.4mm}\begin{minipage}{0.92\linewidth}
      \PanelTitle{内容说明}
      \scriptsize \begin{itemize}
  \item 用 shadow task、Monitor verification、GPC/GPT protection 和 accelerator-specific cleanup，把真实。
\end{itemize}
    \end{minipage}
    \vspace{1mm}
  \end{minipage}}
  \end{column}
  \begin{column}{0.49\textwidth}
    \fcolorbox{Rule}{Panel}{%
  \begin{minipage}[t]{0.97\linewidth}
    \vspace{1.1mm}
    \hspace{1.4mm}\begin{minipage}{0.92\linewidth}
      \PanelTitle{核心设计拆解}
      \scriptsize \begin{itemize}
  \item 01. Shadow task and Monitor-side verification
  \item 02. GPC/GPT protected accelerator buffers
  \item 03. GPU/FPGA workflowow-specific cleanup
\end{itemize}
    \end{minipage}
    \vspace{1mm}
  \end{minipage}}
  \end{column}
\end{columns}
\SourceFootnote{来源：CAGE confidential accelerator workflowow for Arm CCA｜证据等级：E1 / 同行评审改进证据}
\end{frame}


\begin{frame}[t,shrink=6]{CAGE confidential accelerator workflowow for Arm CCA：实验结果}
\DeckKicker{实验结果}
\SlideMeta{证据等级}{E1 / 同行评审改进证据}{来源状态：本地 PDF 已核验}
{\large\bfseries\color{Ink} 论文报告 GPU benchmark 开销 0.58\%-5.31\%，FPGA benchmark 开销 9.61\%-16.30\%，TCB 增量约。\par}\vspace{1.2mm}
\begin{columns}[T,totalwidth=\textwidth]
  \begin{column}{0.45\textwidth}
    \fcolorbox{Rule}{Paper}{%
  \begin{minipage}[t]{0.97\linewidth}
    \vspace{1.1mm}
    \hspace{1.4mm}\begin{minipage}{0.92\linewidth}
      \PanelTitle{内容说明}
      \scriptsize \begin{itemize}
  \item 论文报告 GPU benchmark 开销 0.58\%-5.31\%，FPGA benchmark 开销 9.61\%-16.30\%，TCB 增量约 1301 LoC + 140 LoC
\end{itemize}
    \end{minipage}
    \vspace{1mm}
  \end{minipage}}
  \end{column}
  \begin{column}{0.49\textwidth}
    \fcolorbox{Rule}{Panel}{%
  \begin{minipage}[t]{0.97\linewidth}
    \vspace{1.1mm}
    \hspace{1.4mm}\begin{minipage}{0.92\linewidth}
      \PanelTitle{实验/证据边界}
      \scriptsize {\tiny
\begin{tabularx}{\linewidth}{@{}YYY@{}}
\toprule
\textbf{对象} & \textbf{可支撑} & \textbf{边界} \\
\midrule
数据/来源 & CAGE confidential accelerator workflowow for Arm CCA & 本地 PDF 已核验 \\
Baseline/对照 & 以论文或规范原文为准 & 不跨材料外推 \\
指标/对象 & 只使用当前来源可证明的信息 & E1 / 同行评审改进证据 \\
使用边界 & 可进入本方向演进图 & 不能替代更强证据 \\
\bottomrule
\end{tabularx}
}
    \end{minipage}
    \vspace{1mm}
  \end{minipage}}
  \end{column}
\end{columns}
\SourceFootnote{来源：CAGE confidential accelerator workflowow for Arm CCA｜证据等级：E1 / 同行评审改进证据}
\end{frame}


\begin{frame}[t,shrink=6]{CAGE confidential accelerator workflowow for Arm CCA：文章评价}
\DeckKicker{文章评价}
\SlideMeta{证据等级}{E1 / 同行评审改进证据}{来源状态：本地 PDF 已核验}
{\large\bfseries\color{Ink} 当前 Arm CCA GPU/FPGA confidential accelerator 的 peer-reviewed SOTA\par}\vspace{1.2mm}
\begin{columns}[T,totalwidth=\textwidth]
  \begin{column}{0.45\textwidth}
    \fcolorbox{Rule}{Paper}{%
  \begin{minipage}[t]{0.97\linewidth}
    \vspace{1.1mm}
    \hspace{1.4mm}\begin{minipage}{0.92\linewidth}
      \PanelTitle{内容说明}
      \scriptsize \begin{itemize}
  \item 设计优势：当前 Arm CCA GPU/FPGA confidential accelerator 的 peer-reviewed SOTA。
  \item 局限性：Device identity、SPDM/TDISP、production CCA hardware、复杂多租户 accelerator scheduling 和侧信道仍。
  \item 商业化潜力：可迁移到 DPU/SmartNIC packet-processing rules、queue descriptors 和 TLS keys；风险在 TCB 增量、设备身。
  \item 证据边界：E1 / Peer-reviewed SOTA；本地 PDF 已核验。
\end{itemize}
    \end{minipage}
    \vspace{1mm}
  \end{minipage}}
  \end{column}
  \begin{column}{0.49\textwidth}
    \fcolorbox{Rule}{Panel}{%
  \begin{minipage}[t]{0.97\linewidth}
    \vspace{1.1mm}
    \hspace{1.4mm}\begin{minipage}{0.92\linewidth}
      \PanelTitle{文章评价}
      \scriptsize {\tiny
\begin{tabularx}{\linewidth}{@{}YY@{}}
\toprule
\textbf{维度} & \textbf{判断} \\
\midrule
设计优势 & 当前 Arm CCA GPU/FPGA confidential accelerator 的 peer-reviewed SOTA。 \\
主要局限 & Device identity、SPDM/TDISP、production CCA hardware、复杂多租户 accelerator... \\
商业落地 & 可迁移到 DPU/SmartNIC packet-processing rules、queue descriptors 和 TLS keys；风险在。 \\
讲稿定位 & 当前 SOTA/补充边界; 证据强度 E1 \\
\bottomrule
\end{tabularx}
}
    \end{minipage}
    \vspace{1mm}
  \end{minipage}}
  \end{column}
\end{columns}
\SourceFootnote{来源：CAGE confidential accelerator workflowow for Arm CCA｜证据等级：E1 / 同行评审改进证据}
\end{frame}


\begin{frame}[t,shrink=6]{Accelerator / DPU / SmartNIC confidential offload：方向总结}
\DeckKicker{方向总结}
\SlideMeta{证据等级}{方向综述}{来源状态：公开材料综合}
{\large\bfseries\color{Ink} Accelerator / DPU / SmartNIC confidential offload 的结论是：三篇主讲材料共同给出技术演进，但仍要用证据等级约束每个结论。\par}\vspace{1.2mm}
\begin{columns}[T,totalwidth=\textwidth]
  \begin{column}{0.45\textwidth}
    \fcolorbox{Rule}{Paper}{%
  \begin{minipage}[t]{0.97\linewidth}
    \vspace{1.1mm}
    \hspace{1.4mm}\begin{minipage}{0.92\linewidth}
      \PanelTitle{内容说明}
      \scriptsize \begin{itemize}
  \item HETEE: rack-scale heterogeneous TEE：开创/基础入口，E1 / 基础证据 peer-reviewed。
  \item CloudScale confidential computing with heterogeneous devices：代表性改进，E1 / 同行评审改进证据。
  \item CAGE confidential accelerator workflowow for Arm CCA：当前 SOTA/补充边界，E1 / 同行评审改进证据。
  \item 本方向结论：先用三篇主讲材料建立演进关系，再用证据等级控制机制、实验和商业化结论。
\end{itemize}
    \end{minipage}
    \vspace{1mm}
  \end{minipage}}
  \end{column}
  \begin{column}{0.49\textwidth}
    \fcolorbox{Rule}{Panel}{%
  \begin{minipage}[t]{0.97\linewidth}
    \vspace{1.1mm}
    \hspace{1.4mm}\begin{minipage}{0.92\linewidth}
      \PanelTitle{本方向方法对比}
      \scriptsize {\tiny
\begin{tabularx}{\linewidth}{@{}YYY@{}}
\toprule
\textbf{论文} & \textbf{定位} & \textbf{选择理由} \\
\midrule
HETEE: rack-scale heterogeneous TEE & 开创/基础入口 & HETEE is a foundational peer-reviewed work for heterogeneous... \\
CloudScale confidential computing with heterogeneous devices & 代表性改进 & Cloud-scale acceleration work is a peer-reviewed 2024 SOTA item for... \\
CAGE confidential accelerator workflowow for Arm CCA & 当前 SOTA/补充边界 & CAGE is a recent peer-reviewed SOTA item for GPU accelerator enclaves. \\
\bottomrule
\end{tabularx}
}
    \end{minipage}
    \vspace{1mm}
  \end{minipage}}
  \end{column}
\end{columns}
\SourceFootnote{来源：论文原文、官方规范或公开材料｜证据等级：见本页 badge}
\end{frame}
